%% ---------------------------------------------------------------------
\section{Para llevar}
\label{sec:vt-cierre}
%% ---------------------------------------------------------------------

\begin{parallevar}
\begin{enumerate}
  \item \textbf{La hora de arranque importa:} para $\dot x=f(t,x)$ la solución depende de $t_0$, y aparecen las nociones uniformes. US: $\delta$ no depende de $t_0$. UAS: tampoco el tiempo de convergencia. En sistemas invariantes en el tiempo (y periódicos) la uniformidad es gratis.
  \item \textbf{Funciones de comparación:} UAS es lo mismo que una sola cota $\norm{x(t)}\le\beta(\norm{x(t_0)},t-t_0)$ con $\beta\in\KL$. Hay que revisar las dos variables de $\beta$: estrictamente creciente en el tamaño, decreciente a cero en el tiempo transcurrido.
  \item \textbf{Lyapunov con $V(t,x)$:} $W_1(x)\le V\le W_2(x)$ y $\Vdot\le-W_3(x)$ dan UAS; con $W_1$ radialmente no acotada, GUAS; con cotas $k_i\norm x^a$, exponencial. ``Decreciente'' ($V\le W_2$) es lo que hace uniforme el resultado, y $\Vdot$ incluye $\partial V/\partial t$.
  \item \textbf{$\Vdot<0$ no basta:} hace falta una cota $-W_3(x)$ independiente de $t$. Tampoco bastan los valores propios congelados (Markus--Yamabe) ni una linealización sin uniformidad.
  \item \textbf{Con $\Vdot$ semidefinida:} LaSalle--Yoshizawa da $W(x(t))\to0$; Barbalat sobre la ecuación de la variable que falta completa la convergencia; en el caso lineal periódico, Floquet decide.
  \item \textbf{Inestabilidad:} es la negación de estabilidad, no ``escape a infinito''. Chetaev (invariante) y los dos teoremas de inestabilidad (variantes) piden una $V$ que sea positiva cerca del origen y crezca donde es positiva; $V$ puede ser indefinida.
\end{enumerate}
\end{parallevar}

\begin{trampa}[title={\faBug\hspace{0.45em}Las trampas del capítulo, en una caja}]{}
\begin{itemize}
  \item Converger no es converger uniformemente (\cref{tr:vt-converge}); un $\delta(\varepsilon,t_0)$ no es un $\delta(\varepsilon)$ (\cref{tr:vt-estable}).
  \item Decaer a cero en $s$ no basta para ser de clase $\KL$ (\cref{tr:vt-kl}).
  \item $\Vdot$ lleva el término $\partial V/\partial t$ (\cref{tr:vt-parcial}); $\Vdot<0$ en cada instante no es $\Vdot\le-W_3(x)$ (\cref{tr:vt-negativa}).
  \item Los valores propios de $A(t)$ no deciden (\cref{tr:vt-congelar}); la linealización necesita un jacobiano uniforme en $t$ (\cref{ej:vt-linealizacion}).
  \item LaSalle no se traslada tal cual a sistemas variantes (\cref{tr:vt-lasalle}).
  \item Inestable no es no acotado (\cref{tr:vt-inestable}); una $V$ de Chetaev se construye, no se adivina (\cref{tr:vt-cetaevx1}).
  \item Un certificado que falla no prueba inestabilidad: con la $g$ rápida del \cref{ej:vt-ganancia} el sistema sigue siendo exponencialmente estable (\cref{prob:vt-grapida}).
\end{itemize}
\end{trampa}

%% ---------------------------------------------------------------------
\section{Problemas}
\label{sec:vt-problemas}
%% ---------------------------------------------------------------------

Los problemas marcados con \etiquetaM{} requieren \matlab/\simulink; los marcados con \etiquetaEstrella{} son de profundización. Cada uno tiene una trampa: búsquela antes de dar la respuesta.

\subsection*{Uniformidad y funciones de comparación}

\begin{problema}[una tasa que se apaga más despacio]\label{prob:vt-raiz}
Para $\dot x=-x/\sqrt{1+t}$: (a) resuelva la ecuación; (b) calcule el tiempo de asentamiento $T(0.1,t_0)$ y evalúelo en $t_0=0$ y $t_0=100$; (c) ¿es el origen UAS?
\end{problema}

\begin{problema}[estable, pero no uniformemente]\label{prob:vt-6tsin}
\etiquetaM{} Para $\dot x=(6t\sin t-2t)\,x$: (a) escriba la solución; (b) con $t_0=2n\pi$, muestre que $x(t_0+\pi)/x(t_0)=e^{(4n+1)(6-\pi)\pi}$; (c) concluya que el origen no es uniformemente estable. (d) Simule desde $x(t_0)=10^{-3}$ con $n=0,1,2$ y grafique en escala logarítmica.
\end{problema}

\begin{problema}[\texorpdfstring{$\KL$}{KL} y flujos]\label{prob:vt-flujo}
Una función $\sigma\in\KL$ generada por $\dot y=-\alpha(y)$ cumple $\sigma(r,0)=r$ y $\sigma(\sigma(r,s_1),s_2)=\sigma(r,s_1+s_2)$. (a) De las funciones de la pausa de la \cref{sec:vt-comparacion}, ¿cuáles son generadas por una ecuación de comparación? (b) Muestre que $re^{-s}$ y $r/(rs+1)$ sí lo son, y recupere $\alpha$ con $\alpha(r)=-\partial\sigma/\partial s\,(r,0)$.
\end{problema}

\subsection*{Teoremas de Lyapunov}

\begin{problema}[decreciente o no {\citep[ejercicio~2.7]{KeBr23}}]\label{prob:vt-decreciente}
Para $V_1=x_1^2(1+\sin^2t)+x_2^2(1+\cos^2t)$ y $V_2=x_1^2+x_2^2(1+t)$, escriba $\alpha_1$ y, si existe, $\alpha_2$. Explique en una frase por qué no existe para $V_2$.
\end{problema}

\begin{problema}[atracción sin exponencial]\label{prob:vt-cubico}
Para $\dot x=-(1+g(t))\,x^3$ con $0\le g(t)\le k$: (a) con $V=x^2/2$ pruebe que el origen es GUAS; (b) pruebe que no es exponencialmente estable. ¿Para qué sirve la cota $g\le k$?
\end{problema}

\begin{problema}[cota cruzada {\citep[ejercicio~4.2]{Ser00}}]\label{prob:vt-seron2}
Para $\dot x_1=-2x_1+g(t)x_2$, $\dot x_2=g(t)x_1-2x_2$ con $\abs{g(t)}\le1$, pruebe que el origen es GES y dé la tasa. ¿Hace falta alguna condición sobre $\dot g$?
\end{problema}

\begin{problema}[rotación y disipación cúbica {\citep[ejercicio~4.3]{Ser00}}]\label{prob:vt-seron3}
Para $\dot x_1=x_2-g(t)x_1\norm x^2$, $\dot x_2=-x_1-g(t)x_2\norm x^2$ con $k\le g(t)\le M$, $k>0$: (a) pruebe que el origen es GUAS; (b) pruebe que no es exponencialmente estable; (c) ¿qué dice la linealización?
\end{problema}

\begin{problema}[la ecuación de Lyapunov diferencial, escalar]\label{prob:vt-lyapdif}
\etiquetaEstrella{} Para $\dot x=-a(t)x$ con $a(t)=2+\sin t$: (a) escriba la ecuación \eqref{eq:vt-lyapdif} con $Q=1$; (b) muestre que $p(t)=\int_t^\infty e^{-2\int_t^\tau a(\sigma)\dif\sigma}\dif\tau$ la resuelve; (c) pruebe que $\tfrac16\le p(t)\le\tfrac12$. \etiquetaM{} (d) Calcule numéricamente el rango de $p$.
\end{problema}

\subsection*{Sistemas lineales variantes y ejemplos}

\begin{problema}[Markus--Yamabe a mano]\label{prob:vt-markus}
Verifique la segunda columna de $\Phi(t,0)$ en el \cref{ej:vt-markus}, compruebe la fórmula de Liouville $\det\Phi(t,0)=\exp\int_0^t\operatorname{tr}A$, y encuentre una condición inicial cuya solución tienda a cero.
\end{problema}

\begin{problema}[la ganancia rápida]\label{prob:vt-grapida}
\etiquetaM{} En el \cref{ej:vt-ganancia} con $g=(1+\sin10t)/2$: (a) verifique que $\dot g\le g$ falla; (b) calcule numéricamente la matriz de monodromía con periodo $\pi/5$ y sus multiplicadores; (c) ¿es estable el origen? ¿Qué dice eso de la $V$ del ejemplo? (d) Explique el módulo de los multiplicadores con la fórmula de Liouville.
\end{problema}

\subsection*{Inestabilidad}

\begin{problema}[¿hasta dónde llega la cota de Chetaev?]\label{prob:vt-cetaevr}
En el \cref{ej:vt-cetaev}, ¿para qué valores de $r$ la cota $\min\{1-r^2/2,\ 1-r\}\norm x^2$ prueba que $\Vdot>0$ en $\Bola r\setminus\{0\}$? ¿Cambia eso la conclusión de inestabilidad?
\end{problema}

\begin{problema}[van der Pol: qué teorema sirve]\label{prob:vt-vdp}
\etiquetaEstrella{} Para van der Pol con $\mu=1$: (a) con $V=\tfrac12\norm x^2$, calcule $\Vdot$ y explique por qué el primer teorema de inestabilidad no aplica directamente; (b) pruebe la inestabilidad con el método indirecto; (c) explique por qué el tiempo de salida de $\Bola1$ desde $(\delta,0)$ es aproximadamente $2\ln(1/\delta)$.
\end{problema}

%% ---------------------------------------------------------------------
\section{Soluciones seleccionadas}
\label{sec:vt-soluciones}
%% ---------------------------------------------------------------------

\paragraph{Problema~\ref{prob:vt-raiz}}
(a) $x(t)=x_0\exp\bigl(-2(\sqrt{1+t}-\sqrt{1+t_0})\bigr)$. (b) $x=0.1x_0$ cuando $\sqrt{1+t}=\sqrt{1+t_0}+\tfrac12\ln10$, es decir $T=\ln10\,\sqrt{1+t_0}+\tfrac14\ln^210$: $T=3.63$ con $t_0=0$ y $T=24.47$ con $t_0=100$. (c) No: $T$ crece como $\sqrt{t_0}$. Es US ($\abs{x(t)}\le\abs{x_0}$) y AS, pero no UAS.

\paragraph{Problema~\ref{prob:vt-6tsin}}
(a) $x(t)=x(t_0)\exp\bigl(6\sin t-6t\cos t-t^2-6\sin t_0+6t_0\cos t_0+t_0^2\bigr)$. (b) Entre $2n\pi$ y $(2n+1)\pi$, $\int6\tau\sin\tau=6(4n+1)\pi$ y $\int2\tau=(4n+1)\pi^2$. (c) Para tener $\abs{x(t)}<\varepsilon$ desde $t_0=2n\pi$ hace falta $\delta\le\varepsilon e^{-(4n+1)(6-\pi)\pi}$, que tiende a cero con $n$: no hay $\delta(\varepsilon)$ uniforme. El exponente total tiende a $-\infty$, así que el origen es AS.

\paragraph{Problema~\ref{prob:vt-flujo}}
(a) Ninguna: (a) y (d) son $\KL$ pero $\sigma(r,0)=r^2$ y $2r$, distintos de $r$; (b) y (c) no son $\KL$. (b) $re^{-s}$ viene de $\dot y=-y$ y $r/(rs+1)$ de $\dot y=-y^2$; $-\partial_s(re^{-s})|_{s=0}=r$ y $-\partial_s\bigl(r/(rs+1)\bigr)|_{s=0}=r^2$.

\paragraph{Problema~\ref{prob:vt-decreciente}}
$V_1$: $\alpha_1(r)=r^2$ y $\alpha_2(r)=2r^2$, porque $1\le1+\sin^2t\le2$. $V_2$: $\alpha_1(r)=r^2$; en $x=(0,1)$, $V_2=1+t$, que no está acotado por ningún $\alpha_2(1)$.

\paragraph{Problema~\ref{prob:vt-cubico}}
(a) $\Vdot=-(1+g)x^4\le-x^4$, con $V$ invariante en el tiempo, definida positiva y radialmente no acotada: GUAS, con $\abs{x(t)}\le\abs{x_0}/\sqrt{1+2x_0^2(t-t_0)}$. (b) Como $\dot x\ge-(1+k)x^3$ para $x>0$, por comparación $x(t)\ge x_0/\sqrt{1+2(1+k)x_0^2(t-t_0)}$, que decae como $1/\sqrt t$. La cota $g\le k$ no se usa para GUAS: solo para negar la estabilidad exponencial.

\paragraph{Problema~\ref{prob:vt-seron2}}
Con $V=\norm x^2$, $\Vdot=-4\norm x^2+4g\,x_1x_2=x\T\bigl[\begin{smallmatrix}-4&2g\\2g&-4\end{smallmatrix}\bigr]x\le-2\norm x^2$, porque los valores propios son $-4\pm2g\le-2$. Con $k_1=k_2=1$, $k_3=2$, $a=2$: $\norm{x(t)}\le\norm{x(t_0)}e^{-(t-t_0)}$. No hace falta nada sobre $\dot g$: $V$ no depende de $t$.

\paragraph{Problema~\ref{prob:vt-seron3}}
(a) Con $V=\norm x^2$ los términos de rotación se cancelan y $\Vdot=-2g\norm x^4\le-2k\norm x^4$: GUAS. (b) Con $g\le M$, $\frac{\dif}{\dif t}\norm x^2\ge-2M\norm x^4$ y $\norm{x(t)}^2\ge\norm{x_0}^2/(1+2M\norm{x_0}^2(t-t_0))$, que decae como $1/t$. (c) Un centro ($\pm j$): no decide.

\paragraph{Problema~\ref{prob:vt-lyapdif}}
(a) $\dot p-2a(t)p=-1$. (b) Derivando la integral respecto a su límite inferior y a $t$ dentro del exponente, $\dot p=-1+2a(t)p$. (c) Como $1\le a\le3$, $e^{-6(\tau-t)}\le e^{-2\int_t^\tau a}\le e^{-2(\tau-t)}$, e integrando, $\tfrac16\le p\le\tfrac12$. (d) $0.167\le p(t)\le0.467$.

\paragraph{Problema~\ref{prob:vt-markus}}
Para $\psi=e^{-t}(s,c)$: $\dot\psi=e^{-t}(c-s,\ -s-c)$ y $A\psi$ da lo mismo. $\det\Phi=e^{t/2}e^{-t}(c^2+s^2)=e^{-t/2}=e^{\int_0^t(-1/2)}$. Desde $x(0)=(0,1)$, $x(t)=e^{-t}(\sin t,\cos t)\to0$: la inestabilidad no es de todas las soluciones, sino de las que tienen componente en la primera columna.

\paragraph{Problema~\ref{prob:vt-grapida}}
(a) En $t=0$, $g=\tfrac12$ y $\dot g=5$. (b) Los multiplicadores son $0.482\pm0.230j$, con módulo $0.533=e^{-\pi/5}$. (c) Sí, globalmente exponencialmente estable, con tasa $1$: la $V$ del ejemplo dejó de ser un certificado, pero el sistema no cambió de naturaleza. (d) $\operatorname{tr}A=-2$ para todo $t$, así que $\det M=e^{-2\pi/5}$, y con multiplicadores complejos conjugados, $\abs\mu=e^{-\pi/5}$.

\paragraph{Problema~\ref{prob:vt-cetaevr}}
Para $r<1$, porque $1-r>0$ y $1-r^2/2>0$. No: la inestabilidad es local y basta un $r$.

\paragraph{Problema~\ref{prob:vt-vdp}}
(a) $\Vdot=(1-x_1^2)x_2^2\ge0$ cerca del origen, pero se anula en todo el eje $x_2=0$: no es definida positiva. (b) El jacobiano en el origen tiene valores propios $\tfrac12\pm j\tfrac{\sqrt3}2$, con parte real positiva. (c) Cerca del origen la solución es aproximadamente la de la linealización, $\norm{x(t)}\approx\delta e^{t/2}$, que vale $1$ en $t=2\ln(1/\delta)$; el término $O(1)$ corrige la parte en que la no linealidad ya pesa.

\respuestaspausas

\begin{parallevar}[Nota de transparencia]
En la redacción de este capítulo y en la escritura y corrección de los códigos que lo acompañan se contó con la asistencia de Claude, un modelo de inteligencia artificial desarrollado por Anthropic. El contenido fue revisado y validado por el autor, quien asume la responsabilidad por lo aquí expuesto.
\end{parallevar}
